(a) A conic is isomorphic to $\mathbb P^1$ by \exref{I.3.1}(c), whose local rings are localizations of $k[t]$ and are integrally closed.

(b) Each standard chart of $Q_1$ is an affine plane, so its local rings are integrally closed. The charts $x\ne0$ and $y\ne0$ of $Q_2$ are also affine planes. On $w\ne0$, its coordinate ring is $k[x,y,z]/(xy-z^2)\cong k[s^2,t^2,st]$. This ring is integrally closed: an element of its fraction field integral over it lies in the integrally closed ring $k[s,t]$. Multiplication by a nonzero polynomial of even total degree parity preserves the even and odd summands of $k[s,t]$, so a polynomial in the fraction field of the even subring has no odd summand. It therefore belongs to $k[s^2,t^2,st]$.

(c) At the cusp, $t=y/x$ is integral since $t^2=x$, but $t\notin k[t^2,t^3]_{(t^2,t^3)}$. Indeed, if $t=a/b$ with $b(0)\ne0$, the coefficient of $t$ in $tb$ is nonzero, whereas $a\in k[t^2,t^3]$ has no such term.

(d) Localizations of an integrally closed domain are integrally closed, by clearing denominators in an integral equation. Conversely, if all local rings are integrally closed, an element of $K(Y)$ integral over $A(Y)$ lies in every $A(Y)_{\mathfrak m}$ and hence in their intersection $A(Y)$.

(e) Let $B$ be the integral closure of $A(Y)$ in $K(Y)$. By (3.9A), $B$ is a finitely generated $k$-domain, hence is the coordinate ring of a normal affine variety $\widetilde Y$. Inclusion gives $\pi:\widetilde Y\to Y$. For a dominant $\varphi:Z\to Y$ with $Z$ normal, the injection $K(Y)\to K(Z)$ sends every element of $B$ to an element integral over each $\mathcal O_{P,Z}$, hence regular at every $P$. Thus it gives $B\to\mathcal O(Z)$ and a morphism $\theta:Z\to\widetilde Y$ with $\pi\theta=\varphi$. Any such lift induces the same map on $B\subseteq K(Y)$, so it is unique.
